\section{What the formalization checks}
\label{sec:formalization}

The proof, formalization, and exposition were developed with AI assistance.
The accompanying development formalizes the spaces, their topological
properties, both exact size calculations, the finite projective
resolutions, the cocycle, the cup-product calculation, the nonvanishing
argument, and the universal lower bounds. The assembled theorem is about
actual spaces and ordinary sheaf cohomology, rather than an auxiliary
quotient defined only to encode the combinatorial obstruction.

There are three points of translation between the exposition and the
code. First, the integral free sheaf on the whole space is identified
with the constant integral object in Mathlib's definition of ordinary
sheaf cohomology. Second, the cup operation is implemented through the
constant-module Yoneda algebra and the proved comparison of section
complexes from \cref{sec:cup}. The code proves the chain-lift formula;
it does not compare this construction with an independently defined
Mathlib sheaf cup-product operation. Third, although the gauge choices
suggest interpreting \(\eta_r\) as the class of a double cover, a
torsor-classification equivalence is not needed and is not formalized.

The code uses Boolean bits and the field with two elements. Its recursive
configuration type indexed by \(r\) has \(r+1\) coordinates, and its
alternating function has \(r\) factors. The final theorem uses the same
positive degree \(r\) as \cref{thm:main}; some inductive helpers use
an index one smaller.

\subsection{Reproducing the verification}

The project pins Lean v4.31.0 and Mathlib v4.31.0, with Mathlib commit
\begin{center}
\small\nolinkurl{fabf563a7c95a166b8d7b6efca11c8b4dc9d911f}.
\end{center}
The transitive dependency commits are recorded in the manifest. From the
project's \texttt{lean} directory, run
\begin{verbatim}
lake exe cache get
lake build
lake env lean AxiomAudit.lean
\end{verbatim}
The aggregate \texttt{Aoki.lean} imports 64 local modules. The recorded
build succeeds, and the axiom audit for the principal statements and
bridges reports only propositional extensionality, classical choice,
and quotient soundness. The sources contain no proof placeholders or
additional mathematical axioms. The output is retained in
\href{../../../lean/verification.txt}{\texttt{lean/verification.txt}}.

\Cref{app:lean} is a source concordance, not an additional mathematical
assumption: each entry identifies definitions and proved declarations
used in the proof above. File links are relative to the distributed
project, and printed line numbers locate the declarations even in a
viewer that does not open Lean source files.
